\documentclass[11pt]{article}
\usepackage[a4paper,margin=2.5cm]{geometry}
\usepackage{amsmath,amssymb,bm}
\usepackage{graphicx}
\usepackage[colorlinks=true,linkcolor=blue,urlcolor=blue]{hyperref}
\usepackage{enumitem}
\usepackage{tcolorbox}
\usepackage{listings}
\usepackage{xcolor}

\lstset{basicstyle=\ttfamily\small, breaklines=true, columns=fullflexible,
        keywordstyle=\color{blue}, commentstyle=\color{gray}, language=Python}

\newtcolorbox{keybox}[1]{colback=blue!5!white,colframe=blue!60!black,title=#1}
\newtcolorbox{warnbox}[1]{colback=red!5!white,colframe=red!60!black,title=#1}

\title{The Normalizing-Flow $f_0$ Method (Test 10):\\ A Step-by-Step Guide}
\author{\texttt{fit\_h\_flow\_f0\_joint\_mle} in \texttt{fitter.py}}
\date{July 2026}

\begin{document}
\maketitle

\begin{keybox}{The idea in one sentence}
We do not know the initial distribution function $f_0(z_0,w_0)$ of the stars, so we
model it with a \emph{normalizing flow} (a flexible, trainable probability density),
and we fit this flow \emph{jointly} with the dark-matter scale height $h$ by maximizing
the likelihood of the observed snapshot, where the two are connected through
back-integration of the orbits and Liouville's theorem.
\end{keybox}

\tableofcontents

%======================================================================
\section{Setup: what we have and what we want}

\textbf{What we observe.} A snapshot of $N$ stars at time $t_{\rm obs}$:
positions and vertical velocities
\[
\{(z_i, w_i)\}_{i=1}^{N}, \qquad z \text{ in pc},\quad w \text{ in km/s}.
\]

\textbf{What we want.} The dark-matter scale height $h$, which enters the vertical
potential through
\[
\Phi_{\rm DM}(z) \;=\; 4\pi G \rho_{\rm DM}\, h^2 \,\ln\cosh\!\left(\frac{z}{h}\right),
\]
added to the fixed baryonic (thin disk, thick disk, gas) $\mathrm{sech}^2$ components.

\textbf{The catch.} The likelihood of the data depends on the potential \emph{and} on
the (unknown) initial distribution function $f_0$ at $t=0$. In earlier tests $f_0$ was
assumed to be a Gaussian with fixed (possibly wrong) parameters. Test~10 asks: what if
we instead let $f_0$ be a completely flexible learned density --- a normalizing flow?

%======================================================================
\section{Step 1: The physics link --- Liouville's theorem}

Stars move in phase space $(z,w)$ under Hamilton's equations,
\[
\dot z = w, \qquad \dot w = -\frac{\partial \Phi(z,t)}{\partial z}.
\]
The collisionless Boltzmann equation says $f$ is \emph{constant along orbits}:
\[
\frac{df}{dt}=0
\quad\Longrightarrow\quad
f_{\rm obs}(z,w) \;=\; f_0\bigl(z_0(z,w),\, w_0(z,w)\bigr),
\]
where $(z_0, w_0)$ is the point you land on if you take the observed star at
$(z,w)$ and integrate its orbit \textbf{backwards} from $t_{\rm obs}$ to $t=0$
through the trial potential.

\begin{keybox}{Why there is no Jacobian from the dynamics}
In general, when you transform a density through a map you must multiply by the
Jacobian determinant of the map. Here the back-integration map
$(z,w)\mapsto(z_0,w_0)$ is a \emph{Hamiltonian (symplectic)} map, and the code
implements it with a leapfrog integrator, which is symplectic too. Symplectic maps
preserve phase-space volume exactly:
\[
\left|\det \frac{\partial(z_0,w_0)}{\partial(z,w)}\right| \;=\; 1 .
\]
Therefore the density of an observed star is \emph{literally} $f_0$ evaluated at its
back-integrated point --- no correction factor. This is why the loss below is so simple.
\end{keybox}

So the log-likelihood of the snapshot, given a potential (parametrized by $h$, plus an
optional neural-network correction $\delta\Phi$) and given $f_0$, is
\begin{equation}
\log \mathcal{L}
= \sum_{i=1}^{N} \log f_0\Bigl( z_0(z_i,w_i;\, h,\delta\Phi),\;
                                 w_0(z_i,w_i;\, h,\delta\Phi) \Bigr).
\label{eq:loglike}
\end{equation}
Everything that follows is about (a) how $f_0$ is represented by a flow, and
(b) how \eqref{eq:loglike} is maximized over the flow \emph{and} $h$ at the same time.

%======================================================================
\section{Step 2: What a normalizing flow is}

A normalizing flow represents an arbitrary density by \textbf{deforming a simple one}.
Take a base random variable $\bm u \in \mathbb{R}^2$ with a standard normal density
\[
p_{\rm base}(\bm u) = \frac{1}{2\pi} \exp\!\left(-\tfrac12 \|\bm u\|^2\right),
\]
and an \emph{invertible, differentiable} map $T$ (the ``bijection'') from base space to
data space, $\bm x = T(\bm u)$. The classical change-of-variables formula gives the
density of $\bm x$:
\begin{equation}
\log p(\bm x)
= \log p_{\rm base}\bigl(T^{-1}(\bm x)\bigr)
+ \log \left| \det \frac{\partial\, T^{-1}(\bm x)}{\partial \bm x} \right| .
\label{eq:cov}
\end{equation}
Read it as: push the data point back to base space, score it under the Gaussian, and
add the log-Jacobian that accounts for how the map stretches or compresses volume.

Two requirements drive the whole architecture:
\begin{enumerate}[label=(\roman*)]
  \item $T^{-1}$ must be cheap to evaluate (we need it for every star, every epoch);
  \item the Jacobian determinant must be cheap (a general $2\times2$ determinant is
        cheap anyway, but the design below makes it \emph{trivial} and generalizes to
        any dimension).
\end{enumerate}

%======================================================================
\section{Step 3: The bijection --- RealNVP affine coupling layers}

The map $T$ is a stack of $K$ \emph{affine coupling layers}
(class \texttt{\_AffineCoupling}), each acting on the 2-D point. One layer works
like this:

\begin{enumerate}[label=\textbf{3.\arabic*}]
\item \textbf{Split.} Write the point as $(a, b)$: coordinate $a$ is the
      \emph{pass-through}, coordinate $b$ is the one being \emph{transformed}.
      A flag \texttt{flip} alternates the roles between layers (even layers transform
      $w$, odd layers transform $z$), so over the stack both coordinates get moved.

\item \textbf{Condition.} A small MLP takes the pass-through coordinate $a$ and
      outputs two numbers: a log-scale $s(a)$ and a shift $t(a)$:
      \[
      \text{MLP: } 1 \to 64 \to 64 \to 2, \quad \tanh \text{ activations.}
      \]
      The log-scale is soft-bounded, $s = 3\tanh(\hat s/3)$, so no single layer can
      stretch by more than $e^{\pm 3}\approx 20\times$ --- a numerical-stability guard.

\item \textbf{Affine transform.} Only $b$ changes:
      \begin{align}
      \text{forward (data $\to$ base, used in \texttt{log\_prob}):}\quad
          & u = \bigl(b - t(a)\bigr)\, e^{-s(a)}, \qquad a \mapsto a, \\[2pt]
      \text{inverse (base $\to$ data, used for sampling):}\quad
          & b = u\, e^{s(a)} + t(a), \qquad\quad\; a \mapsto a.
      \end{align}
\end{enumerate}

\begin{keybox}{Why this is invertible even though it contains a neural network}
You never need to invert the MLP. Given the output, the pass-through coordinate $a$
is right there unchanged, so you can recompute exactly the same $s(a)$ and $t(a)$ and
undo the affine map in closed form. The network only \emph{parametrizes} the affine
map; the map itself is trivially invertible.
\end{keybox}

\textbf{Identity initialization.} The last linear layer of each MLP is initialized to
zeros, so at the start $s=t=0$ and every coupling layer is the identity. The flow
therefore \emph{begins} as an exact standard Gaussian and gradually deforms it during
training. (This mirrors the zero-init convention used elsewhere in \texttt{fitter.py}
--- and recall the ``dead-init'' bug: here only the \emph{last} layer is zero, the
hidden weights are random, so gradients do flow.)

\textbf{Standardization.} Before entering the stack, physical coordinates are
whitened with a \emph{fixed} (non-trained) affine map,
\[
\bm x = \left(\frac{z-\mu_z}{\sigma_z},\; \frac{w-\mu_w}{\sigma_w}\right),
\]
where $\mu,\sigma$ are the mean/std of the cloud back-integrated once at the initial
guess $h_{\rm init}$. This just puts the data at order unity for the network.

%======================================================================
\section{Step 4: The Jacobian --- why it costs nothing}

Take one coupling layer in the forward direction, mapping $(a,b) \to (a,u)$.
Its Jacobian matrix is
\[
\frac{\partial(a,u)}{\partial(a,b)} =
\begin{pmatrix}
1 & 0 \\[4pt]
\dfrac{\partial u}{\partial a} & e^{-s(a)}
\end{pmatrix}.
\]
The matrix is \textbf{triangular}: $u$ depends on both inputs, but $a$ depends only on
itself. The determinant of a triangular matrix is the product of its diagonal, so the
messy off-diagonal term (which involves derivatives of the neural network!) is
irrelevant:
\begin{equation}
\boxed{\;\log\left|\det \frac{\partial(a,u)}{\partial(a,b)}\right| = -\,s(a)\;}
\end{equation}
That single number per layer is exactly what \texttt{forward} returns alongside the
transformed point. For the stack of $K$ layers, log-determinants simply add:
\[
\log|\det J_{\rm total}| = \sum_{k=1}^{K} \bigl(-s_k\bigr)
\;\;-\;\; \log(\sigma_z \sigma_w),
\]
where the last term is the (constant) Jacobian of the whitening step --- it converts
the density back to \emph{physical units}, per $\mathrm{pc}\cdot\mathrm{km/s}$, so it
is directly comparable to a physical $f_0$.

Putting Steps 2--4 together, \texttt{ZWNormalizingFlow.log\_prob} computes:
\begin{equation}
\log f_0^{\rm flow}(z_0,w_0)
= \underbrace{-\tfrac12\|\bm u\|^2 - \log 2\pi}_{\text{Gaussian base}}
\;+\; \underbrace{\sum_{k=1}^{K} (-s_k)}_{\text{coupling Jacobians}}
\;-\; \underbrace{\log(\sigma_z\sigma_w)}_{\text{whitening}} .
\label{eq:flowlogprob}
\end{equation}

%======================================================================
\section{Step 5: The flow \emph{as} $f_0$ --- assembling the loss}

Here is the conceptual switch of Test 9/10 relative to Test 4:

\begin{center}
\begin{tabular}{l|l|l}
 & \textbf{Test 4} & \textbf{Tests 9/10} \\ \hline
flow models\dots & the \emph{observed} snapshot $f_{\rm obs}$ & the \emph{initial} DF $f_0$ \\
$f_0$ is\dots & known (double Gaussian) & unknown, learned \\
fit matches\dots & flow$(z,w)$ to $f_0(z_0,w_0)$ & data to flow via \eqref{eq:loglike} \\
\end{tabular}
\end{center}

Substituting the flow \eqref{eq:flowlogprob} into the Liouville likelihood
\eqref{eq:loglike}, the training loss for one epoch is
\begin{equation}
\mathcal{L}(\theta_{\rm flow},\, h,\, \theta_{\rm NN}, A_{\rm NN})
= -\frac{1}{|B|}\sum_{i\in B}
   \log f_0^{\rm flow}\Bigl(
      z_0\bigl(z_i,w_i;\, h, A_{\rm NN}\,\delta\Phi_{\theta_{\rm NN}}\bigr),\;
      w_0\bigl(\cdots\bigr)
   \Bigr),
\label{eq:loss}
\end{equation}
where $B$ is a random mini-batch of 500 stars. Note carefully what depends on what:
\begin{itemize}
\item $h$ (and the optional NN potential) sit \emph{inside the arguments}: they change
      \emph{where} each star lands after back-integration;
\item the flow parameters $\theta_{\rm flow}$ change \emph{how probable} any landing
      point is.
\end{itemize}
The back-integration \texttt{torch\_back\_integrate} is written entirely in PyTorch
(1000 leapfrog steps), so automatic differentiation propagates gradients \emph{through
the whole orbit} into $h$: the optimizer literally feels ``if $h$ were slightly
larger, star $i$ would land at a slightly different $(z_0,w_0)$, where the flow
assigns a different probability.''

\textbf{The optional NN potential.} When enabled, a time-dependent correction
$A_{\rm NN}\,\delta\Phi(z,t)$ is added to the force during back-integration. Its
amplitude is hard-capped by projection so that its peak never exceeds 10\% of the
static potential's peak (the \texttt{nn\_frac\_bound} cap): after every optimizer step,
the peak of $|A_{\rm NN}\delta\Phi|$ on a $(z,t)$ grid is measured and $A_{\rm NN}$ is
rescaled down if it exceeds the budget.

%======================================================================
\section{Step 6: Joint minimization --- one optimizer, four parameter groups}

A single Adam optimizer updates everything simultaneously, with per-group learning
rates tuned to each parameter's natural scale:

\begin{center}
\begin{tabular}{l l l}
\textbf{parameters} & \textbf{lr} & \textbf{role} \\ \hline
flow weights $\theta_{\rm flow}$ & $10^{-3}$ & reshape $f_0$ \\
$h$ (a single scalar \texttt{nn.Parameter}) & $2.0$ & move the potential \\
NN potential weights $\theta_{\rm NN}$ & $10^{-2}$ & shape of $\delta\Phi(z,t)$ \\
amplitude $A_{\rm NN}$ & $5\times10^{-2}$ & size of $\delta\Phi$ \\
\end{tabular}
\end{center}

One epoch (of 200) does:
\begin{enumerate}[label=\arabic*.]
\item draw a random subsample of 500 observed stars;
\item back-integrate them $t_{\rm obs}\!\to\!0$ with the \emph{current} $h$ and
      $\delta\Phi$ (differentiable leapfrog);
\item evaluate the loss \eqref{eq:loss} with the \emph{current} flow;
\item backpropagate; clip the gradient on $h$ (norm 50) and on the NN (norm 1);
\item Adam step on all groups at once;
\item clamp $h$ into its bounds and project $A_{\rm NN}$ onto the 10\% cap.
\end{enumerate}

So the flow is chasing whatever cloud the current $h$ produces, while $h$ is
simultaneously being pulled toward whatever makes the cloud most probable under the
current flow.

%======================================================================
\section{Step 7: Why this fails --- the Liouville degeneracy (the point of Test 10)}

\begin{warnbox}{The trap}
The same Liouville theorem that makes the loss elegant makes it \emph{flat in $h$}.
For (almost) \emph{any} trial $h$, the back-integration produces \emph{some} smooth
cloud at $t=0$ --- and a sufficiently flexible flow can fit \emph{that} cloud just as
well as any other. Wrong-$h$ clouds are just as ``fittable'' as the right one. The
likelihood then barely prefers the true $h$, and nothing anchors it except the flow's
\emph{limited capacity}.
\end{warnbox}

Formally: if $\Psi_h$ denotes the back-map at scale height $h$, then for any $h$ the
choice $f_0 = f_{\rm obs}\circ \Psi_h^{-1}$ (the pushforward of the data) achieves the
same maximal likelihood. A rigid parametric $f_0$ (a Gaussian) cannot follow this
family, so it \emph{breaks} the degeneracy --- that is where the constraint on $h$
actually comes from.

Test 10 measures this directly with the \textbf{init-spread}
$|h(\text{init }180) - h(\text{init }350)|$: run the same fit from two starting
guesses; if the likelihood really pins $h$, both must converge to the same value.
The pre-run results (N = data size, flow size = couplings $\times$ hidden):

\begin{center}
\begin{tabular}{l c c}
\textbf{model of $f_0$} & \textbf{init-spread, $N{=}2000$} & \textbf{init-spread, $N{=}500$} \\ \hline
rigid Gaussian (even with wrong $\sigma_w$) & $\sim 1$ pc & $\sim 1$ pc \\
small flow ($8\times64$) & $\sim 3$ pc & $\sim 12$ pc \\
big flow ($16\times128$) & $\sim 8$ pc & $\sim 23$ pc \\
\end{tabular}
\end{center}

More flexibility or less data $\Rightarrow$ weaker constraint, exactly as the
degeneracy argument predicts. And the NN potential does not rescue the rigid-but-biased
Gaussian either: from one init it stays idle, from the other it \emph{destabilizes}
$h$ (246 $\to$ 297).

\begin{keybox}{Takeaway}
The information about $h$ in a snapshot-likelihood fit comes from \emph{restricting}
$f_0$, not from the data alone. Relaxing $f_0$ into a flow destroys the constraint,
and adding potential flexibility (the NN) cannot rebuild it. This dead-end is what
motivates changing the \emph{objective}: instead of a snapshot likelihood, use a
statistic that directly targets the non-equilibrium term, e.g.\ a collisionless-Boltzmann
residual $\partial f/\partial t$.
\end{keybox}

%======================================================================
\appendix
\section{Appendix: mapping the math to the code}

\begin{center}
\begin{tabular}{l l}
\textbf{math object} & \textbf{code (\texttt{fitter.py})} \\ \hline
coupling bijection, $\pm s$ log-det & \texttt{\_AffineCoupling.forward / .inverse} \\
soft bound $s = 3\tanh(\hat s /3)$ & \texttt{\_scale\_shift} \\
change of variables \eqref{eq:flowlogprob} & \texttt{ZWNormalizingFlow.log\_prob} \\
whitening $\mu,\sigma$ & \texttt{set\_standardization} (from $h_{\rm init}$ cloud) \\
back-map $\Psi_h$ (symplectic leapfrog) & \texttt{torch\_back\_integrate} \\
joint loss \eqref{eq:loss} & \texttt{loss = -flow.log\_prob(z0, w0).mean()} \\
joint fit driver & \texttt{fit\_h\_flow\_f0\_joint\_mle} \\
10\% cap projection & \texttt{\_vpeak()} + \texttt{A\_nn.mul\_(A\_nn\_max/vp)} \\
Test 10 experiment & notebook cells \texttt{test10-md}, \texttt{test10-code} \\
\end{tabular}
\end{center}

\end{document}
